
\documentclass{article}
\usepackage{colortbl}
\usepackage{makecell}
\usepackage{multirow}
\usepackage{supertabular}

\begin{document}

\newcounter{utterance}

\centering \large Interaction Transcript for game `bbh', experiment `salient\_translation\_error\_detection', episode 23 with reflexion{-}Qwen3.5{-}9B.
\vspace{24pt}

{ \footnotesize  \setcounter{utterance}{1}
\setlength{\tabcolsep}{0pt}
\begin{supertabular}{c@{$\;$}|p{.15\linewidth}@{}p{.15\linewidth}p{.15\linewidth}p{.15\linewidth}p{.15\linewidth}p{.15\linewidth}}
    \# & \multicolumn{2}{c}{Player} && \multicolumn{2}{c}{Game Master} \\
    \hline

    \theutterance \stepcounter{utterance}  
    & & & \multicolumn{4}{p{0.6\linewidth}}{
        \cellcolor[rgb]{0.9,0.9,0.9}{
            \makecell[{{p{\linewidth}}}]{
                \texttt{\tiny{[P1$\langle$GM]}}
                \texttt{Detect the type of error in an English translation of a German source sentence.} \\
\\ 
\texttt{Q: The following translations from German to English contain a particular error. That error will be one of the following types: Named Entities: An entity (names, places, locations, etc.) is changed to a different entity. Numerical Values: Numerical values (ordinals or cardinals), dates, and/or units are changed. Modifiers or Adjectives: The modifiers and adjectives pertaining to a noun are changed. Negation or Antonyms: Introduce or remove a negation or change comparatives to their antonyms. Facts: Trivial factual errors not pertaining to the above classes are introduced in the translations. Dropped Content: A significant clause in the translation is removed. Please identify that error.  Source: In der Liste der Baudenkmale in Lenzen (Elbe) sind alle Baudenkmale der brandenburgischen Stadt Lenzen (Elbe) und ihrer Ortsteile aufgelistet.} \\
\texttt{Translation: In the list of architectural monuments in Lenzen all architectural monuments of the Brandenburg city of Lenzen and its districts are listed.} \\
\texttt{The translation contains an error pertaining to} \\
\texttt{Options:} \\
\texttt{(A) Modifiers or Adjectives} \\
\texttt{(B) Numerical Values} \\
\texttt{(C) Negation or Antonyms} \\
\texttt{(D) Named Entities} \\
\texttt{(E) Dropped Content} \\
\texttt{(F) Facts} \\
\texttt{A: (D)} \\
\\ 
\texttt{Q: The following translations from German to English contain a particular error. That error will be one of the following types: Named Entities: An entity (names, places, locations, etc.) is changed to a different entity. Numerical Values: Numerical values (ordinals or cardinals), dates, and/or units are changed. Modifiers or Adjectives: The modifiers and adjectives pertaining to a noun are changed. Negation or Antonyms: Introduce or remove a negation or change comparatives to their antonyms. Facts: Trivial factual errors not pertaining to the above classes are introduced in the translations. Dropped Content: A significant clause in the translation is removed. Please identify that error.  Source: Auf dieser Seite sind die Baudenkmäler der oberbayerischen Großen Kreisstadt Landsberg am Lech zusammengestellt.} \\
\texttt{Translation: On this page are compiled the architectural monuments of the town of Landsberg am Lech.} \\
\texttt{The translation contains an error pertaining to} \\
\texttt{Options:} \\
\texttt{(A) Modifiers or Adjectives} \\
\texttt{(B) Numerical Values} \\
\texttt{(C) Negation or Antonyms} \\
\texttt{(D) Named Entities} \\
\texttt{(E) Dropped Content} \\
\texttt{(F) Facts} \\
\texttt{A: (E)} \\
\\ 
\texttt{Q: The following translations from German to English contain a particular error. That error will be one of the following types: Named Entities: An entity (names, places, locations, etc.) is changed to a different entity. Numerical Values: Numerical values (ordinals or cardinals), dates, and/or units are changed. Modifiers or Adjectives: The modifiers and adjectives pertaining to a noun are changed. Negation or Antonyms: Introduce or remove a negation or change comparatives to their antonyms. Facts: Trivial factual errors not pertaining to the above classes are introduced in the translations. Dropped Content: A significant clause in the translation is removed. Please identify that error.  Source: Łeba ist eine Kleinstadt und ein Badeort im Powiat Lęborski der polnischen Woiwodschaft Pommern.} \\
\texttt{Translation: Eba is not a small town and seaside resort in the Powiat Léborski county of the Pomeranian Voivodeship of Poland.} \\
\texttt{The translation contains an error pertaining to} \\
\texttt{Options:} \\
\texttt{(A) Modifiers or Adjectives} \\
\texttt{(B) Numerical Values} \\
\texttt{(C) Negation or Antonyms} \\
\texttt{(D) Named Entities} \\
\texttt{(E) Dropped Content} \\
\texttt{(F) Facts} \\
\texttt{A: (C)} \\
\\ 
\texttt{Q: The following translations from German to English contain a particular error. That error will be one of the following types: Named Entities: An entity (names, places, locations, etc.) is changed to a different entity. Numerical Values: Numerical values (ordinals or cardinals), dates, and/or units are changed. Modifiers or Adjectives: The modifiers and adjectives pertaining to a noun are changed. Negation or Antonyms: Introduce or remove a negation or change comparatives to their antonyms. Facts: Trivial factual errors not pertaining to the above classes are introduced in the translations. Dropped Content: A significant clause in the translation is removed. Please identify that error.  Source: Richard Raphael Roland Risse war ein deutscher Historien{-}, Genre{-} und Bildnismaler der Düsseldorfer Schule.} \\
\texttt{Translation: Risse was a German historical, genre and portrait painter of the Düsseldorf School.} \\
\texttt{The translation contains an error pertaining to} \\
\texttt{Options:} \\
\texttt{(A) Modifiers or Adjectives} \\
\texttt{(B) Numerical Values} \\
\texttt{(C) Negation or Antonyms} \\
\texttt{(D) Named Entities} \\
\texttt{(E) Dropped Content} \\
\texttt{(F) Facts} \\
\texttt{A:} \\
            }
        }
    }
     \\ \\

    \theutterance \stepcounter{utterance}  
    & \multicolumn{4}{p{0.6\linewidth}}{
        \cellcolor[rgb]{0.561,0.737,0.561}{
            \makecell[{{p{\linewidth}}}]{
                \texttt{\tiny{[P1$\rangle$GM]}}
                \texttt{(D)} \\
\\ 
\texttt{**Analysis:**} \\
\texttt{The source text contains the full name "Richard Raphael Roland Risse". The translation shortens this to just "Risse". While this is a common practice in English to avoid redundancy, strictly speaking, the full named entity (the person's full name) present in the source has been altered/omitted in the translation. In the context of the provided error categories, changing a specific named entity (the full name) to a different form (just the surname) fits best under **Named Entities**, as the specific identity string provided in the source is not fully preserved.} \\
\\ 
\texttt{However, let's re{-}evaluate based on the strict definitions often used in these datasets (like WMT or similar translation error detection tasks):} \\
\texttt{{-} **Named Entities**: An entity is changed to a *different* entity. (e.g., "Berlin" {-}> "London").} \\
\texttt{{-} **Dropped Content**: A significant clause is removed.} \\
\texttt{{-} **Facts**: Trivial factual errors.} \\
\\ 
\texttt{Is shortening a full name to a surname considered changing it to a *different* entity? Usually, no. It's often considered a normalization or a drop of modifiers. But looking at the previous examples:} \\
\texttt{1. "brandenburgischen Stadt" {-}> "city of" (Dropped modifier/adjective? Or just rephrasing). The error identified was (D) Named Entities? Wait, in the first example: "brandenburgischen Stadt Lenzen" vs "city of Lenzen". The error identified in the prompt's example answer was (D). Why? Perhaps because "brandenburgischen" (Brandenburg) was dropped, which is a location modifier, but the answer key says (D). Let's look closer at Example 1.} \\
\texttt{   {-} Source: "...brandenburgischen Stadt Lenzen..."} \\
\texttt{   {-} Translation: "...city of Lenzen."} \\
\texttt{   {-} The answer provided in the prompt for Example 1 is **(D)**. This suggests that dropping the specific regional descriptor ("Brandenburg") from a place name is classified as a **Named Entities** error in this specific schema, or perhaps the schema considers the full proper noun phrase "Brandenburg city of Lenzen" as the entity and the translation "city of Lenzen" as a different entity representation.} \\
\\ 
\texttt{Let's look at Example 2:} \\
\texttt{{-} Source: "...oberbayerischen Großen Kreisstadt Landsberg am Lech..."} \\
\texttt{{-} Translation: "...town of Landsberg am Lech."} \\
\texttt{{-} Answer: **(E) Dropped Content**. Here, "oberbayerischen Großen Kreisstadt" (Upper Bavarian large district town) was dropped. This is clearly dropped content.} \\
\\ 
\texttt{Why was Example 1 (D) and Example 2 (E)?} \\
\texttt{In Example 1, "brandenburgischen" is an adjective modifying the city. In Example 2, "oberbayerischen Großen Kreisstadt" is a long descriptive phrase.} \\
            }
        }
    }
    & & \\ \\

    \theutterance \stepcounter{utterance}  
    & & & \multicolumn{2}{p{0.3\linewidth}}{
        \cellcolor[rgb]{0.9,0.9,0.9}{
            \makecell[{{p{\linewidth}}}]{
                \texttt{\tiny{[GM$|$GM]}}
                \texttt{Target: (D)} \\
            }
        }
    }
    & & \\ \\

    \theutterance \stepcounter{utterance}  
    & & & \multicolumn{2}{p{0.3\linewidth}}{
        \cellcolor[rgb]{0.9,0.9,0.9}{
            \makecell[{{p{\linewidth}}}]{
                \texttt{\tiny{[GM$|$GM]}}
                \texttt{game\_result = LOSE} \\
            }
        }
    }
    & & \\ \\

    \theutterance \stepcounter{utterance}  
    & & \multicolumn{4}{p{0.6\linewidth}}{
        \cellcolor[rgb]{0.95,0.95,0.95}{
            \makecell[{{p{\linewidth}}}]{
                \texttt{\tiny{[P1$\leftrightarrow$P]}}
                \texttt{Post{-}episode Reflexion (not sent to GM):} \\
\texttt{New plan: I incorrectly classified the shortening of a full name ("Richard Raphael Roland Risse" to "Risse") as a Named Entities error because I focused on the change in the string itself. However, based on the pattern in the examples (where dropping a regional descriptor like "brandenburgischen" was labeled Dropped Content in the second example, and the first example's logic seems to hinge on whether the core entity identity changed vs. descriptive details being omitted), shortening a name to its surname is a removal of descriptive modifiers, not a change to a different entity. Therefore, I should have selected **(E) Dropped Content** for the third example, as the specific names "Richard Raphael Roland" were omitted. My new strategy is to strictly distinguish between "changing the entity to a different one" (e.g., Berlin {-}> Paris) and "omitting descriptive parts of the entity" (e.g., Full Name {-}> Surname, or City + Region {-}> City). I will re{-}evaluate the third example as (E) and apply this distinction to future cases.} \\
            }
        }
    }
    & \\ \\

\end{supertabular}
}

\end{document}
